%% ma: L10-B15-0109
%% lop: 10
%% chuong: 6
%% bai: 15
%% dang: 10.15.7
%% loai: TLN
%% muc: TH
%% dap_an: "2"
%% nguon: "Ảnh giáo viên gửi 25/11/2026 (lô 2): Bài tập trắc nghiệm – Toán 10 Bài 15 Hàm số (sự biến thiên, đồ thị), Phần III Câu 1"
%% kien_thuc: "Lớp 10 Bài 15 (giá trị của hàm số cho bởi nhiều công thức)"
%% trang_thai: cho-duyet
%% ly_do: "Dạng 10.15.7 mới đề xuất."
\begin{ex}
Cho $f(x)=\begin{cases}x-4 & \text{khi } x\ge0\\ x^2-4x+1 & \text{khi } x<0\end{cases}$. Có bao nhiêu giá trị nguyên của tham số $m$ để $f(m^2)+f(-2)=18$?
\shortans{2}
\loigiai{$f(-2)=4+8+1=13$; $m^2\ge0$ nên $f(m^2)=m^2-4$. Phương trình: $m^2-4+13=18\Leftrightarrow m^2=9\Leftrightarrow m=\pm3$. Có $2$ giá trị.}
\end{ex}
